\section{Results}
\label{sec:results}

\begin{table}[!t]
\caption{Lifetime, lossy channel, 30 paired runs (mean $\pm$ sample std, rounds). FND/HND/LND are first, half and last node death. AUC is the area under the living-node curve, the censoring-robust survival measure.}
\label{tab:lifetime}
\centering
\begin{tabular}{llrrrr}
\toprule
Protocol & Gen. & FND & HND & LND & AUC ($\times 10^3$) \\
\midrule
Direct (no clustering) & -- & 114 $\pm$ 11 & 1,144 $\pm$ 183 & 3,202 $\pm$ 78 & 132.0 \\
LEACH & G1 & 1,046 $\pm$ 43 & 1,679 $\pm$ 25 & 2,191 $\pm$ 81 & 167.1 \\
PEGASIS & G1 & 1,588 $\pm$ 331 & 2,199 $\pm$ 9 & 2,630 $\pm$ 358 & 216.4 \\
TEEN & G1 & 2,208 $\pm$ 164 & 5,601 $\pm$ 189 & $>$6,999\textsuperscript{c} & 535.4 \\
APTEEN & G1 & 2,144 $\pm$ 139 & 5,367 $\pm$ 161 & $>$6,999\textsuperscript{c} & 518.4 \\
NSGA-II & G2 & 1,873 $\pm$ 31 & 1,965 $\pm$ 23 & 2,697 $\pm$ 256 & 200.6 \\
Fuzzy T2 & G2 & 1,462 $\pm$ 50 & 2,006 $\pm$ 28 & 2,109 $\pm$ 48 & 191.9 \\
SOM & G3 & 959 $\pm$ 127 & 1,998 $\pm$ 32 & 2,874 $\pm$ 115 & 196.8 \\
DQN & G3 & 1,472 $\pm$ 55 & 1,920 $\pm$ 31 & 2,325 $\pm$ 59 & 192.1 \\
GCN & G3 & 335 $\pm$ 142 & 1,824 $\pm$ 70 & 3,185 $\pm$ 142 & 180.2 \\
\bottomrule
\end{tabular}

{\footnotesize \textsuperscript{c} Censored: the network had not died at the 7000-round cap in any run. Use AUC, not LND, for these two.\par}
\end{table}

\begin{table}[!t]
\caption{Delivery, lossy channel. Packets at the base station measure aggregation depth, not delivery, and are not comparable across protocols; readings delivered and data yield are. Data yield is readings arriving at the sink divided by readings taken, so it is the column that prices TEEN and APTEEN's suppression.}
\label{tab:delivery}
\centering
\begin{tabular}{llrrrr}
\toprule
Protocol & Gen. & Packets at BS & Readings & PDR (\%) & Data yield \\
 & & ($\times 10^3$) & delivered ($\times 10^3$) & & \\
\midrule
Direct (no clustering) & -- & 128.1 & 128.1 & 96.9 & 0.970 \\
LEACH & G1 & 11.0 & 146.6 & 87.7 & 0.877 \\
PEGASIS & G1 & 2.3 & 182.2 & 84.2 & 0.842 \\
TEEN & G1 & 21.2 & 58.6 & 90.8 & 0.109 \\
APTEEN & G1 & 21.5 & 62.6 & 90.8 & 0.121 \\
NSGA-II & G2 & 9.7 & 179.9 & 89.6 & 0.896 \\
Fuzzy T2 & G2 & 9.1 & 175.4 & 91.4 & 0.914 \\
SOM & G3 & 9.4 & 175.8 & 89.3 & 0.893 \\
DQN & G3 & 9.0 & 173.5 & 90.3 & 0.903 \\
GCN & G3 & 8.6 & 162.8 & 90.3 & 0.903 \\
\bottomrule
\end{tabular}
\end{table}


\subsection{What Clustering Actually Buys}

The direct-transmission baseline makes the first result concrete. LEACH reaches
first node death at 1{,}046 rounds against the baseline's 114, an improvement of
9.2 times. But the baseline's last node dies at 3{,}202 rounds against LEACH's
2{,}191, which is 46\% later. Clustering does not extend lifetime; it
redistributes it, buying a much later first death at the cost of an earlier last
one.

\begin{figure}[!t]
\centering
\includegraphics[width=\columnwidth]{fig1_alive_nodes.png}
\caption{Living nodes against round, lossy channel, mean over 30 paired trials.
The direct baseline loses its first node almost immediately and then decays
slowly, while the clustered protocols hold a full population far longer and then
collapse. The two shapes cross, which is why no single death point orders these
protocols and why the area under this curve is reported alongside all three.}
\label{fig:alive}
\end{figure}

The mechanism is geometric. With the sink at 150~m and a crossover at 87.7~m, 36
of the 100 nodes in the seed-0 topology lie close enough to transmit on the
free-space branch and are individually cheap to operate; the expected fraction
over the placement distribution is 32.7\%. Under direct transmission those nodes
are never asked to relay anyone else's traffic and simply outlive everything
LEACH produces. Under rotation they subsidize the expensive far-field nodes.
Fig.~\ref{fig:alive} shows the two survival shapes crossing.

This is why every lifetime claim in this paper is reported at three death points
plus the area under the living-node curve. Any single figure silently picks a
winner, and the two most commonly reported figures pick opposite ones.

The baseline also produces the clearest illustration of survivorship bias in
delivery statistics. Its delivery ratio of 96.9\% is the highest of any
configuration on the lossy channel, well above LEACH's 87.7\%, and it is an
artifact: the distant nodes that lose the most packets die first, so the
surviving population is progressively closer to the sink. Measured directly, the
baseline's delivery ratio rises from 83.7\% over its first hundred rounds to
98.6\% over its last hundred. Delivery ratio is never reported here without the
living-node curve beside it.

\subsection{Lifetime Across the Three Generations}

\begin{figure}[!t]
\centering
\includegraphics[width=\columnwidth]{figB_fnd_distribution.png}
\caption{First node death per trial, lossy channel, all 30 paired trials shown
individually rather than as a mean and a deviation. The reactive pair leads, the
GCN is visibly bimodal, and PEGASIS carries one topology at 449 rounds against a
body near 1{,}800. A mean and standard deviation would describe none of these
three distributions correctly.}
\label{fig:fnd-dist}
\end{figure}

Table~\ref{tab:lifetime} gives the full lossy-channel results. Against LEACH,
with paired permutation testing and Holm correction across the family, every
protocol differs significantly. The reactive pair leads on first death by a wide
margin, TEEN at 2{,}208 rounds with $\Delta = +1{,}163$, followed by NSGA-II at
1{,}873 ($\Delta = +827$), PEGASIS at 1{,}588, the DQN at 1{,}472 and the fuzzy
system at 1{,}462. Two protocols are significantly worse than LEACH: the
self-organizing map at 959 and the GCN at 335.

Three observations survive scrutiny and are worth stating separately from the
ranking.

NSGA-II posts both the best non-reactive first death and the tightest spread of
any of the nine clustering protocols, with a standard deviation of 31 rounds
against LEACH's 43, the DQN's 55 and PEGASIS's 331. Its advantage holds in both
channel configurations. But its convergence is modest, and we report that as a
finding rather than suppress it. Over the logged mid-run reclustering, the
population's best round-energy objective improves by 1.76\% and its best
minimum-head-energy objective by 3.61\%, while only cluster-size balance moves
substantially, from 2.99 to 1.85 for the population best and from 9.19 to 5.58
at the knee point. With five heads drawn from a hundred candidates, random
initialization already lands near a good solution. NSGA-II's advantage here
should therefore not be attributed to the sophistication of the search.

PEGASIS has by far the largest first-death variance of any protocol, 331 rounds
against a typical 30 to 60. A greedy chain is sensitive to where the topology
happens to place its outliers, and a single badly positioned node changes the
chain's whole cost structure; one trial gives 449 rounds against a body near
1{,}800. Its half-death figure, by contrast, has a standard deviation of 9
rounds, the tightest in the study, because once the chain has been rebuilt a few
times its structure becomes topology-insensitive.

The GCN's first death is not merely poor but bimodal. Across the first five
seeds it records 448, 394, 418, 112 and 98 rounds, and over all 30 trials eight
fall between 94 and 151 while the remaining twenty-two fall between 284 and 482.
A mean of 335 with a standard deviation of 142 describes neither mode.
Fig.~\ref{fig:fnd-dist} shows the shape, and
Section~\ref{sec:credit} explains it.

\subsection{A Conclusion the Channel Sweep Retracts}
\label{sec:retraction}

\begin{table*}[!t]
\caption{First-node-death difference against LEACH, in rounds, under both channel configurations. Paired sign-flip permutation test (20\,000 permutations) with paired-bootstrap confidence intervals, Holm-corrected across the nine comparisons within each channel. \textbf{Fuzzy T2 and DQN are significant under loss and not significant without it}, which is the study's central negative result.}
\label{tab:paired-fnd}
\centering
\begin{tabular}{lrcrrcr}
\toprule
& \multicolumn{3}{c}{Lossy channel} & \multicolumn{3}{c}{Ideal channel} \\
\cmidrule(lr){2-4}\cmidrule(lr){5-7}
Protocol & $\Delta$FND & 95\% CI & $p_{\mathrm{Holm}}$ & $\Delta$FND & 95\% CI & $p_{\mathrm{Holm}}$ \\
\midrule
Direct (no clustering) & -932 & [-944, -919] & 0.00045 & -1,153 & [-1,162, -1,144] & 0.00045 \\
PEGASIS & +542 & [+418, +647] & 0.00045 & +188 & [+78, +293] & 0.009 \\
TEEN & +1,163 & [+1,109, +1,217] & 0.00045 & +2,593 & [+2,542, +2,647] & 0.00045 \\
APTEEN & +1,099 & [+1,055, +1,143] & 0.00045 & +2,434 & [+2,392, +2,480] & 0.00045 \\
NSGA-II & +827 & [+811, +844] & 0.00045 & +482 & [+471, +493] & 0.00045 \\
Fuzzy T2 & +417 & [+393, +440] & 0.00045 & +17 & [+0, +34] & 0.119 \\
SOM & -87 & [-135, -41] & 0.001 & -399 & [-443, -357] & 0.00045 \\
DQN & +426 & [+404, +450] & 0.00045 & +22 & [-4, +49] & 0.124 \\
GCN & -710 & [-764, -660] & 0.00045 & -1,140 & [-1,193, -1,090] & 0.00045 \\
\bottomrule
\end{tabular}
\end{table*}

\begin{table}[!t]
\caption{Where the energy goes and where the heads sit, measured by the engine's own accounting categories (seed 0, 1100 rounds). Packet error is near zero below 120~m, so the \emph{mean} head distance cannot explain the retry gap --- every mean is below that. The tail can: across the six single-hop protocols, the share of head-rounds beyond 120~m predicts retry energy share with $r = 0.95$. PEGASIS is the exception that proves the rule: its leader rotates uniformly and 30\% of leader-rounds sit beyond 120~m, yet it retries barely at all, because only the leader ever faces the sink.}
\label{tab:energy-split}
\centering
\begin{tabular}{lrrrrr}
\toprule
Protocol & Retry & Control & \multicolumn{2}{c}{Head-to-sink (m)} & Heads $>$120 m \\
\cmidrule(lr){4-5}
 & (\%) & (\%) & mean & p90 & (\% of rounds) \\
\midrule
LEACH & 6.39 & 12.73 & 102.1 & 143.5 & 29.9 \\
PEGASIS & 1.39 & 4.38 & 102.2 & 143.5 & 30.0 \\
SOM & 5.28 & 2.38 & 102.8 & 136.4 & 31.1 \\
NSGA-II & 3.57 & 2.46 & 94.2 & 127.5 & 17.7 \\
Fuzzy T2 & 3.65 & 2.33 & 85.6 & 122.1 & 11.5 \\
DQN & 2.31 & 3.95 & 81.1 & 111.2 & 5.1 \\
GCN & 2.25 & 3.85 & 77.1 & 104.6 & 3.2 \\
\bottomrule
\end{tabular}
\end{table}


Applying the decision rule of Section~\ref{sec:channel-pair} to the 30-run
sweep, one conclusion does not survive, and it is a load-bearing one. On the
lossy channel the fuzzy system and the DQN improve first death over LEACH by 417
and 426 rounds respectively, both at the corrected significance floor of
$p_{\mathrm{Holm}} = 0.00045$. On the ideal channel the same comparisons give 17
and 22 rounds at $p_{\mathrm{Holm}} = 0.119$ and $0.124$, which is not
significant. Remove packet loss and both protocols become statistically
indistinguishable from LEACH on the metric this literature most often reports.
Table~\ref{tab:paired-fnd} gives both channels side by side.

We measured the mechanism rather than arguing it, and the measurement rules out
the explanation we expected. Table~\ref{tab:energy-split} splits each protocol's
energy by accounting category over a single seed and reports the full
distribution of head-to-sink distance in the same run, not only its mean.

\begin{figure}[!t]
\centering
\includegraphics[width=\columnwidth]{figF_head_distance.png}
\caption{Distribution of head-to-sink distance per protocol, drawn over the
packet-error curve. Every protocol's \emph{mean} falls in the flat region where
packet error is effectively zero, so the mean cannot be what drives the retry
gap. The distributions differ in their upper tails, and only the tails reach the
waterfall.}
\label{fig:head-dist}
\end{figure}

The obvious reading of these numbers would be that LEACH's average cluster head
sits at 102.1~m, past the 87.7~m crossover and therefore onto the fourth-power
branch, while the DQN's sits at 81.1~m. That is true, and it explains part of
the difference in raw transmit energy, but it cannot explain the
retransmissions. Packet error is $2.6\times10^{-8}$ at 100~m and does not reach
$10^{-3}$ until 120~m, and every protocol in Table~\ref{tab:energy-split} has a
mean head distance below 120~m. A protocol whose heads all sat at the mean would
retry essentially never.

What separates them is the upper tail, shown in Fig.~\ref{fig:head-dist}. LEACH
spends 29.9\% of its head-rounds beyond 120~m and 6.39\% of its budget on
retransmissions. The DQN spends 5.1\% of its head-rounds beyond 120~m and 2.31\%
on retransmissions. Across the six single-hop protocols the share of head-rounds
beyond 120~m correlates with retry energy share at $r = 0.95$. The mean distance
correlates almost as strongly, at $r = 0.93$, because the two quantities are
collinear; the reason to prefer the tail is not a better fit but that only the
tail lies in the region where packets actually fail.

\begin{figure}[!t]
\centering
\includegraphics[width=\columnwidth]{figC_mechanism.png}
\caption{Retry energy share against the share of head-rounds beyond 120~m, lossy
channel, seed 0. The six single-hop protocols line up at $r = 0.95$. PEGASIS,
plotted separately, has the same 30\% tail as LEACH and almost no retry cost,
because only its leader ever faces the sink.}
\label{fig:mechanism}
\end{figure}

PEGASIS is the exception that fixes the scope of the claim, and we state it
rather than exclude it quietly. Its leader rotates uniformly over chain
positions, so its leader-distance distribution is identical to the node
distribution and 30.0\% of leader-rounds sit beyond 120~m, the same as LEACH.
Yet it retries on only 1.39\% of its budget. Including it drops the correlation
from 0.95 to 0.50. The reason is structural: in a chain, exactly one node faces
the sink per round while every other transmission is a short hop to a neighbor,
so the tail of the sink-facing distribution is paid once per round instead of
five times. The claim is therefore specific to single-hop clustered protocols,
where every head pays the sink hop, and Fig.~\ref{fig:mechanism} plots it that
way.

The honest claim is therefore narrower, and more interesting, than ``the
soft-computing and learned protocols cluster better''. Their demonstrable skill
is keeping the far tail of the head-to-sink hop out of the error waterfall, and
that skill pays only in proportion to how lossy that hop is. Delete the loss by
making the channel lossless and the advantage disappears. The same table shows
LEACH's compensating strength: it consumes 20\% more total energy than the DQN
yet matches its first death on a clean channel, because its epoch rule rotates
head duty perfectly.

The retraction is specific to first death and we do not overstate it. On area
under the living-node curve and on readings delivered, every protocol's
difference from LEACH is significant at the corrected floor in both channels, in
the same direction and at similar magnitude; the fuzzy system gains 24{,}793
node-rounds on the lossy channel and 22{,}120 on the ideal.
Table~\ref{tab:paired-robust} gives these. The correct statement is that the
advantage in overall survival and in data delivered is robust to the channel,
and only the advantage in first death is not. First death is a single order
statistic and the most fragile of the four lifetime measures.

\begin{table}[!t]
\caption{Survival and delivery differences against LEACH under both channels. Unlike first node death (Table~\ref{tab:paired-fnd}), every one of these gaps holds in both configurations, so these are the conclusions the study can state without a channel caveat.}
\label{tab:paired-robust}
\centering
\begin{tabular}{lrrrr}
\toprule
& \multicolumn{2}{c}{Lossy} & \multicolumn{2}{c}{Ideal} \\
\cmidrule(lr){2-3}\cmidrule(lr){4-5}
Protocol & $\Delta$ ($\times 10^3$) & $p_{\mathrm{Holm}}$ & $\Delta$ ($\times 10^3$) & $p_{\mathrm{Holm}}$ \\
\midrule
\multicolumn{5}{l}{\textit{AUC}} \\
\quad Direct (no clustering) & -35.1 & 0.0004 & -32.8 & 0.0004 \\
\quad PEGASIS & +49.2 & 0.0004 & +44.5 & 0.0004 \\
\quad TEEN & +368.3 & 0.0004 & +389.6 & 0.0004 \\
\quad APTEEN & +351.3 & 0.0004 & +371.7 & 0.0004 \\
\quad NSGA-II & +33.5 & 0.0004 & +31.7 & 0.0004 \\
\quad Fuzzy T2 & +24.8 & 0.0004 & +22.1 & 0.0004 \\
\quad SOM & +29.7 & 0.0004 & +29.2 & 0.0004 \\
\quad DQN & +25.0 & 0.0004 & +22.5 & 0.0004 \\
\quad GCN & +13.1 & 0.0004 & +11.1 & 0.0004 \\
\addlinespace
\multicolumn{5}{l}{\textit{Readings}} \\
\quad Direct (no clustering) & -18.5 & 0.0004 & -31.8 & 0.0004 \\
\quad PEGASIS & +35.6 & 0.0004 & +44.2 & 0.0004 \\
\quad TEEN & -88.1 & 0.0004 & -106.0 & 0.0004 \\
\quad APTEEN & -84.0 & 0.0004 & -101.3 & 0.0004 \\
\quad NSGA-II & +33.2 & 0.0004 & +32.2 & 0.0004 \\
\quad Fuzzy T2 & +28.8 & 0.0004 & +22.4 & 0.0004 \\
\quad SOM & +29.2 & 0.0004 & +29.3 & 0.0004 \\
\quad DQN & +26.9 & 0.0004 & +22.5 & 0.0004 \\
\quad GCN & +16.2 & 0.0004 & +9.8 & 0.0004 \\
\bottomrule
\end{tabular}
\end{table}

\begin{table}[!t]
\caption{Every ranking that moves when the channel changes: 4 of 40 protocol--metric ranks. All are single-place swaps between adjacent protocols. Ordering is therefore stable; it is the \emph{significance} of specific first-node-death gaps that is not (Table~\ref{tab:paired-fnd}).}
\label{tab:robustness}
\centering
\begin{tabular}{llrrr}
\toprule
Metric & Protocol & Lossy rank & Ideal rank & Shift \\
\midrule
hnd & Fuzzy T2 & 4 & 5 & +1 \\
hnd & SOM & 5 & 4 & -1 \\
readings\_delivered\_total & Fuzzy T2 & 4 & 5 & +1 \\
readings\_delivered\_total & DQN & 5 & 4 & -1 \\
\bottomrule
\end{tabular}
\end{table}


Across all forty protocol-metric rankings, four ranks move between channel
configurations, and both moving pairs are near-ties between adjacent protocols
(Table~\ref{tab:robustness}). The ordering is therefore stable. What is not
stable is the statistical significance of specific first-death gaps.

\subsection{What Temporal Credit Assignment Buys}
\label{sec:credit}

\begin{table}[!t]
\caption{Head-duty concentration over 1100 rounds (seed 0), in rounds served as cluster head per node. LEACH's 55 is exact rather than approximate: its epoch mechanism elects every node exactly once per 20 rounds. GCN's busiest node serves 13.6 times as often, which is the direct cause of its early first node death (Table~\ref{tab:lifetime}) --- it does not choose badly, it chooses the same good node until that node dies.}
\label{tab:head-rotation}
\centering
\begin{tabular}{lrrrr}
\toprule
Protocol & Busiest node & Median & Never served & Gini \\
\midrule
LEACH & 55 & 55 & 0 & 0.001 \\
PEGASIS & 11 & 11 & 0 & 0.000 \\
NSGA-II & 95 & 60 & 0 & 0.207 \\
Fuzzy T2 & 360 & 40 & 4 & 0.416 \\
SOM & 175 & 50 & 10 & 0.393 \\
DQN & 223 & 38 & 16 & 0.476 \\
GCN & 750 & 16 & 13 & 0.700 \\
\bottomrule
\end{tabular}
\end{table}


The DQN and the GCN consume identical inputs and differ only in architecture and
training signal, so the gap between them isolates one variable. Over the 30
paired trials the DQN's first death ranges from 1{,}361 to 1{,}580, a spread of
219 rounds with a standard deviation of 55. The GCN's ranges from 94 to 482, a
spread of 388 with a standard deviation of 142, and it is bimodal rather than
merely wide.

\begin{figure}[!t]
\centering
\includegraphics[width=\columnwidth]{figG_head_rotation.png}
\caption{Rounds served as cluster head per node over 1{,}100 rounds, seed 0,
sorted. LEACH's epoch rule produces a flat line at exactly 55. The GCN's curve is
the opposite extreme: it concentrates 750 rounds of head duty on one node while
13 nodes never serve at all, which is what its single-round objective is
indifferent to.}
\label{fig:rotation}
\end{figure}

Table~\ref{tab:head-rotation} measures why over 1{,}100 rounds on seed 0. The
GCN's busiest node serves as cluster head for 750 of those rounds against
LEACH's 55, a factor of 13.6, and its Gini coefficient over head duty is 0.700
against LEACH's 0.001. Thirteen of its nodes never serve at all, and its median
node serves 16 rounds. It does not choose badly. It repeatedly chooses the same
genuinely good node until that node dies, which is exactly the failure a
single-round objective cannot see. The DQN, with the same four features
including rounds since last serving as head, concentrates far less: 223 rounds
on its busiest node and a Gini of 0.476. Fig.~\ref{fig:rotation} shows the two
curves against LEACH's flat line.

This is the paper's cleanest controlled contrast. The GCN can see rotation state
in its input. What it lacks is any term in its loss that rewards deferring cost
to extend lifetime, and the cost of that omission is a first death 3.1 times
earlier than LEACH's. It is also the only protocol to lose to LEACH on first
death in both channel configurations by a wide margin.

\subsection{Delivery and the Reactive Tradeoff}

TEEN and APTEEN dominate every lifetime metric and fail to reach full depletion
within the 7{,}000-round cap in all 30 trials, so their last-death figures are
lower bounds. Their area under the living-node curve is 3.2 times LEACH's. This
is not free. Their data yield is 0.109 and 0.121 against roughly 0.9 for the
proactive protocols, and they deliver 58{,}551 and 62{,}579 readings against
LEACH's 146{,}616. They survive by not reporting, and their reporting rate,
0.15 and 0.17 of rounds respectively, is the quantity that makes the lifetime
figure interpretable. Reported without it, their lifetime is not a result about
clustering at all.

Table~\ref{tab:delivery} also shows why packets at the base station must not be
compared across protocols. PEGASIS delivers 182{,}179 readings in 2{,}272
packets, an aggregation ratio of 80, while the direct baseline delivers 128{,}075
readings in 128{,}075 packets. On a packet count PEGASIS looks like the worst
protocol in the study and on readings it is the best.

\subsection{Cost and Latency}

\begin{table*}[!t]
\caption{Cost, lossy channel. Counted operations are the honest cross-protocol comparison; wall-clock setup time reflects implementation quality as much as algorithm design and is reported separately for that reason. DQN and GCN setup times are inference only: their one-off offline training is not included here (Section~\ref{sec:validity}).}
\label{tab:cost}
\centering
\begin{tabular}{llrrrrrr}
\toprule
Protocol & Gen. & Latency & p95 & Control pkts & Ops per & Setup & Peak mem \\
 & & (ms) & (ms) & ($\times 10^3$) & round & (ms) & (MB) \\
\midrule
Direct (no clustering) & -- & 600 & 1,488 & 0.0 & 0 & 0.00 & 6.0 \\
LEACH & G1 & 274 & 837 & 171.6 & 329 & 0.82 & 1.7 \\
PEGASIS & G1 & 1,045 & 1,626 & 213.9 & 59 & 0.05 & 1.5 \\
TEEN & G1 & 419 & 990 & 573.9 & 307 & 0.83 & 5.1 \\
APTEEN & G1 & 407 & 979 & 554.3 & 296 & 0.80 & 5.1 \\
NSGA-II & G2 & 387 & 573 & 10.8 & 86,348 & 102.89 & 2.1 \\
Fuzzy T2 & G2 & 401 & 1,028 & 9.7 & 2,301 & 5.14 & 1.6 \\
SOM & G3 & 388 & 538 & 10.7 & 6,194 & 38.28 & 2.0 \\
DQN & G3 & 402 & 918 & 11.9 & 8,210 & 0.52 & 1.7 \\
GCN & G3 & 392 & 981 & 12.2 & 5,454 & 0.88 & 2.3 \\
\bottomrule
\end{tabular}
\end{table*}


Latency separates the protocols structurally. The clustered protocols sit
between 387 and 419~ms mean, LEACH is lowest at 274~ms because its stochastic
election produces smaller average clusters, and PEGASIS is an outlier at
1{,}045~ms mean and 1{,}626~ms at the 95th percentile. Chain depth is the
dominant latency cost and no amount of energy efficiency compensates for it. The
ordering is structurally correct, but the absolute figures are lower bounds,
since the model contains no queueing or contention delay.

Computational cost is where the three measurement notions of
Section~\ref{sec:metrics} disagree most sharply, which is the argument for
keeping them separate. NSGA-II costs 102.9~ms per round of setup against LEACH's
0.82~ms, a factor of 126, and 86{,}348 counted operations against 329, a factor
of 262. The self-organizing map is second at 38.3~ms. But the ordering by
counted operations and by wall-clock is not the same. The DQN performs 8{,}210
counted operations per round in 0.52~ms, while the fuzzy system performs 2{,}301
in 5.14~ms. The DQN does 3.6 times the algorithmic work in a tenth of the time,
because its work is a dense matrix product and the fuzzy system's is an
iterative type-reduction loop. Reporting either number alone would invert the
comparison.

Control traffic differs by more than an order of magnitude across classes:
573{,}870 packets for TEEN and 171{,}644 for LEACH against 9{,}657 to 12{,}216
for the centralized five. That gap is substantially an artifact of the
accounting model rather than of clustering quality, and
Section~\ref{sec:validity} quantifies it.

\subsection{Node Count, Field Area, and the Regime Boundary}
\label{sec:scale}

\begin{table*}[!t]
\caption{First node death (rounds) across node counts and field areas, lossy channel, 15 paired runs per cell.}
\label{tab:scale-fnd}
\centering
\begin{tabular}{llrrrrrrrrr}
\toprule
  &   & \multicolumn{3}{c}{50$\times$50 m} & \multicolumn{3}{c}{100$\times$100 m} & \multicolumn{3}{c}{150$\times$150 m} \\
\cmidrule(lr){3-5}\cmidrule(lr){6-8}\cmidrule(lr){9-11}
Protocol & Gen. & $N{=}50$ & $N{=}100$ & $N{=}150$ & $N{=}50$ & $N{=}100$ & $N{=}150$ & $N{=}50$ & $N{=}100$ & $N{=}150$ \\
\midrule
Direct (no clustering) & -- & 2,178 & 1,943 & 1,677 & 120 & 112 & 111 & 23 & 22 & 22 \\
LEACH & G1 & 1,791 & 1,617 & 1,586 & 641 & 1,034 & 1,085 & 155 & 303 & 354 \\
PEGASIS & G1 & 2,097 & 2,114 & 2,062 & 1,550 & 1,610 & 1,517 & 663 & 815 & 720 \\
TEEN & G1 & 6,808 & 5,899 & 5,186 & 2,028 & 2,201 & 2,064 & 425 & 550 & 532 \\
APTEEN & G1 & 6,831 & 5,684 & 5,002 & 1,947 & 2,141 & 1,992 & 427 & 519 & 527 \\
NSGA-II & G2 & 2,184 & 2,173 & 2,160 & 1,818 & 1,872 & 1,865 & 717 & 1,052 & 1,077 \\
Fuzzy T2 & G2 & 1,656 & 1,685 & 1,698 & 1,379 & 1,454 & 1,470 & 381 & 393 & 577 \\
SOM & G3 & 1,382 & 1,280 & 1,296 & 1,153 & 965 & 994 & 591 & 342 & 328 \\
DQN & G3 & 1,500 & 1,628 & 1,677 & 1,296 & 1,478 & 1,537 & 454 & 559 & 833 \\
GCN & G3 & 200 & 208 & 138 & 152 & 371 & 305 & 71 & 158 & 82 \\
\bottomrule
\end{tabular}

{\footnotesize The base station scales with the field at $(W/2,\,1.5W)$, so field size sets channel severity as well as area. Node-to-sink distance by field size --- 50$\times$50~m: mean 52~m, 0\% of nodes beyond 120~m; 100$\times$100~m: mean 104~m, 33\% of nodes beyond 120~m; 150$\times$150~m: mean 156~m, 75\% of nodes beyond 120~m. At $50\times50$~m no link reaches the error waterfall at all, which is why the advantage of the optimized and learned protocols over LEACH collapses there.\par}
\end{table*}

\begin{table}[!t]
\caption{Cells ordered by node density. If density drove lifetime the rows would trend with it; they do not --- the numbers track the mean sink distance column instead. Averaged over all ten protocols, tripling the node count changes first node death by a factor of 0.88--1.26, while tripling the field side changes it by a factor of 5--7. The two cells closest in density (50 and 44 nodes/ha) differ by 1.7$\times$ for NSGA-II, 2.1$\times$ for LEACH and 5.5$\times$ for direct transmission. Note also that \emph{clustering only pays once the sink is far}: in the $50\times50$~m rows the no-clustering baseline outlives LEACH.}
\label{tab:scale-density}
\centering
\begin{tabular}{llrrrrr}
\toprule
$N$ & Field (m) & Density & Mean sink & \multicolumn{3}{c}{First node death (rounds)} \\
\cmidrule(lr){5-7}
 & & (nodes/ha) & dist.\ (m) & Direct & LEACH & NSGA-II \\
\midrule
50 & 150$\times$150 & 22 & 156 & 23 & 155 & 717 \\
100 & 150$\times$150 & 44 & 156 & 22 & 303 & 1,052 \\
50 & 100$\times$100 & 50 & 104 & 120 & 641 & 1,818 \\
150 & 150$\times$150 & 67 & 156 & 22 & 354 & 1,077 \\
100 & 100$\times$100 & 100 & 104 & 112 & 1,034 & 1,872 \\
150 & 100$\times$100 & 150 & 104 & 111 & 1,085 & 1,865 \\
50 & 50$\times$50 & 200 & 52 & 2,178 & 1,791 & 2,184 \\
100 & 50$\times$50 & 400 & 52 & 1,943 & 1,617 & 2,173 \\
150 & 50$\times$50 & 600 & 52 & 1,677 & 1,586 & 2,160 \\
\bottomrule
\end{tabular}
\end{table}


The headline study measures one deployment. The $3\times3$ sweep of
Section~\ref{sec:scale-design} asks which of its conclusions are about
clustering and which are about that deployment.
Table~\ref{tab:scale-fnd} gives first node death in all nine cells and
Table~\ref{tab:scale-density} reorders the same cells by node density.

\begin{figure}[!t]
\centering
\includegraphics[width=\columnwidth]{figD_scale_fnd.png}
\caption{First node death across the nine cells, log scale, color by
generation. Moving between field sizes shifts every protocol by roughly an order
of magnitude; moving between node counts within a field size barely moves
anything. The reactive pair reverses rank between the smallest and largest
fields.}
\label{fig:scale-fnd}
\end{figure}

\emph{Area dominates node count by roughly five to one.} Averaged over all ten
configurations, tripling the node count changes first node death by a factor of
0.88 to 1.26, while tripling the field side changes it by a factor of 5 to 7.
Fig.~\ref{fig:scale-fnd} shows the separation directly.

\emph{The driver is distance to the sink, not density.}
Table~\ref{tab:scale-density} orders the cells by node density, and if density
were the cause the rows would trend with it. They do not; they track the mean
sink distance column instead. The two cells closest in density, 50 and 44
nodes per hectare, differ by 1.7 times for NSGA-II, 2.1 times for LEACH and 5.5
times for direct transmission, because one has its sink 104~m away and the other
156~m. We note the limitation this rests on: the grid was not designed to hold
density constant, so no two cells match exactly, and this is a near-match rather
than a controlled contrast. The effect sizes differ by enough that the
conclusion is not sensitive to the gap, but the design does not license a
stronger statement.

\begin{figure}[!t]
\centering
\includegraphics[width=\columnwidth]{figE_scale_density.png}
\caption{First node death against node density, with points coloured by field
size. Density spans a factor of 27 across the grid and explains almost none of
the variation; the three field-size bands separate cleanly instead.}
\label{fig:scale-density}
\end{figure}

\emph{Clustering is harmful when the sink is close.} In the $50\times50$~m
field, where the mean node-to-sink distance is 52~m and no node lies beyond
120~m, LEACH's first node death relative to no clustering at all is a factor of
0.82, 0.83 and 0.95 at $N = 50$, 100 and 150. The no-clustering baseline outlives
LEACH in all three cells, on area under the living-node curve as well as on first
death. At $100\times100$~m the same ratios are 5.4, 9.3 and 9.8, and at
$150\times150$~m they are 6.7, 13.8 and 16.1. The opening result of
Section~\ref{sec:results} is therefore not a statement about clustering as such
but about an operating regime: clustering pays once the sink hop is expensive
enough to be worth aggregating away, and not before.

\emph{The channel retraction reproduces from an independent direction.} In the
$50\times50$~m field the DQN's first-death advantage over LEACH falls to $+10$
rounds at $N = 100$ and to $-291$ at $N = 50$, against $+444$ at
$100\times100$~m; the fuzzy system falls to $+68$ and $-135$ against $+419$.
This is the same boundary found in Section~\ref{sec:retraction} by removing
packet loss, reached instead by moving the sink close enough that no link
reaches the waterfall. The fuzzy system is what makes the attribution safe: it
has no training of any kind, so its collapse cannot be a transfer failure and
must be a property of the environment. Since the DQN behaves the same way, the
DQN's collapse is a property of the environment too.

\emph{Rankings are partly stable and the exceptions are informative.} NSGA-II
holds ranks 1 to 3 in every cell, takes rank 1 in all three $150\times150$~m
cells, and beats LEACH in all nine by margins from $+394$ to $+1{,}177$ rounds.
The GCN holds ranks 9 to 10 everywhere. Against those, the DQN's rank improves
monotonically as the field grows, and TEEN and APTEEN reverse outright, from
ranks 1 and 2 in the two smaller fields to ranks 4 to 6 at $150\times150$~m
where NSGA-II and PEGASIS overtake them. Rankings by area under the living-node
curve are considerably more stable than by first death, with every protocol
except the direct baseline moving at most four places.

\emph{The frozen policies do not collapse off their training distribution.}
The DQN and the GCN carry $N = 100$, $100\times100$~m weights into all nine
cells with no retraining. Neither fails catastrophically, and the GCN's deficit
against LEACH actually shrinks as the field grows, from between $-1{,}409$ and
$-1{,}591$ rounds in the smallest field to between $-84$ and $-272$ in the
largest, although it remains worse than LEACH in every cell. This is weaker
evidence than a retrained comparison would give, and it is not a claim that
these policies are scale-invariant. It is a claim that their reported advantage
is not an artifact of evaluating them only at the point where they were trained.

Two of our five pre-registered expectations about this sweep were wrong, and we
report both. We expected the generational gap to widen at $150\times150$~m; it
narrows there as well as at $50\times50$~m, so the gap is largest in the middle
cell. We expected node density to matter; it does not.
